\ifdefined\mainfile\else
\documentclass{article}
\usepackage{graphicx}
\usepackage{float}
\usepackage{amsmath}
\usepackage{amssymb}
\usepackage[a4paper, left=2.5cm, right=2.5cm, top=2.5cm, bottom=2.5cm]{geometry}
\usepackage[colorlinks=true,linkcolor=black,citecolor=blue,urlcolor=blue]{hyperref}
\begin{document}
\fi

\section{Program Counter}
\label{sec:program-counter}
% 8-bit Program Counter med PS(1:0) styring
% PS=00: Hold, PS=01: Increment, PS=10: Branch (PC + Offset), PS=11: Jump

% \begin{figure}[H]
%     \centering
%     \includegraphics[width=0.75\linewidth]{images/pc/TODO_pc_diagram.png}
%     \caption{Program Counter blokdiagram}
%     \label{fig:pc-overview}
% \end{figure}

\subsection{Design}
\label{sec:pc-design}
% Beskrivelse af PC logik og PS kontrolsignaler
husk at referere 
\begin{figure}[H]
    \centering
    \includegraphics[width=1\linewidth]{image.png}
    \caption{Program counter 4-bit}
    \label{fig:placeholder}
\end{figure}
\subsection{VHDL Implementering}
\label{sec:pc-vhdl}
% VHDL kode og forklaring

\subsection{Test}
\label{sec:pc-test}
% Testbench og verifikation -- Hold, Increment, Branch, Jump

\ifdefined\mainfile\else
\end{document}
\fi
